% Palimpsest — default LaTeX header for `pal pdf` (pandoc + XeLaTeX).
% Book-like prose typography plus multilingual glyph fallback (Greek, Hebrew,
% arrows, letterlike aleph, CJK). Uses only packages available in BasicTeX.
% Fonts below ship with macOS; override main/mono per book via book.toml [pdf]
% (main_font / mono_font), or point [pdf] header at your own file entirely.

\usepackage{fancyhdr}

% --- Prose typography (book-like) ---------------------------------------
\usepackage{microtype}                 % subtle kerning/justification polish
\usepackage{setspace}\doublespacing     % double-spaced manuscript format (12pt, 1in margins)
\setlength{\parindent}{0pt}            % block paragraphs (match the reading copy — no indent)
\setlength{\parskip}{0.7\baselineskip plus 2pt}  % a visible space between paragraphs
\raggedbottom                          % don't stretch short pages to full height
\frenchspacing                         % single space after sentences

% --- Fonts ---------------------------------------------------------------
% Baskerville: a transitional serif. Menlo for the code/notation passages.
% Both ship with macOS. Override via book.toml [pdf] main_font / mono_font.
\setmainfont{Baskerville}[
  BoldFont={Baskerville Bold},
  ItalicFont={Baskerville Italic},
  BoldItalicFont={Baskerville Bold Italic},
]
\setmonofont{Menlo}[Scale=0.85]

% --- Multilingual glyph fallback ----------------------------------------
% The main serif lacks Hebrew and some Greek/arrow glyphs. Auto-switch fonts per
% Unicode block via XeTeX interchar classes. Fonts named below ship with macOS.
\XeTeXinterchartokenstate=1
\newcount\ebloop

% Generic helper: assign all codepoints in [#2,#3) to char-class #1.
\def\ebrange#1#2#3{%
  \ebloop=#2
  \loop \ifnum\ebloop<#3
    \XeTeXcharclass\ebloop=#1
    \advance\ebloop by 1
  \repeat}
\def\ebbind#1#2{% open with font #2 on entering class #1, close on leaving
  \XeTeXinterchartoks 0 #1 = {\bgroup#2}%
  \XeTeXinterchartoks 255 #1 = {\bgroup#2}%
  \XeTeXinterchartoks 4095 #1 = {\bgroup#2}%
  \XeTeXinterchartoks #1 0 = {\egroup}%
  \XeTeXinterchartoks #1 255 = {\egroup}%
  \XeTeXinterchartoks #1 4095 = {\egroup}}

% Greek (U+0370–03FF) and Greek Extended (U+1F00–1FFF).
\newfontfamily\greekfont{Times New Roman}
\newXeTeXintercharclass\greekclass
\ebrange\greekclass{"0370}{"0400}
\ebrange\greekclass{"1F00}{"2000}
\ebbind\greekclass\greekfont

% Hebrew (U+0590–05FF).
\newfontfamily\hebrewfont{Arial Hebrew}
\newXeTeXintercharclass\hebrewclass
\ebrange\hebrewclass{"0590}{"0600}
\ebbind\hebrewclass\hebrewfont

% Arrows (U+2190–21FF) + letterlike aleph (transfinite orders).
\newfontfamily\arrowfont{Apple Symbols}
\newXeTeXintercharclass\arrowclass
\ebrange\arrowclass{"2190}{"2200}
\ebrange\arrowclass{"2135}{"2139}
\ebbind\arrowclass\arrowfont

% CJK Unified Ideographs (U+4E00–9FFF) — defensive.
\newfontfamily\cjkfont{Songti SC}
\newXeTeXintercharclass\cjkclass
\ebrange\cjkclass{"4E00}{"A000}
\ebbind\cjkclass\cjkfont
